\section{The original Cech differential is a derivation}

\begin{lemma}[The whole ordered Cech product rule]
\label{lem:native-cech-ordered-cup-leibniz}
\uses{def:native-cech-ordered-cup}
\leanok
\lean{EnvelopeProperty.DI.CechAlexanderWhitneyProduct.cup_cech_leibniz}
For the original relative Frobenius forms and the original all-tuple
sheaf Cech complex, the actual ordered product satisfies the whole
sheaf-map identity
\[
 \delta\circ\smile=\smile\circ(\delta\otimes1)
      +(-1)^j\smile\circ(1\otimes\delta),
 \qquad \check C^j(E^q)\otimes\check C^{j'}(E^r)
       \longrightarrow\check C^{j+j'+1}(E^{q+r}).
\]
The two output indices $(j+1)+j'$ and $j+(j'+1)$ are transported
by their original equalities with $j+j'+1$. These are the original
indexed objects and original cup morphisms. No covering, finiteness,
or smoothness assertion is needed for this identity.
\end{lemma}
\begin{proof}
\leanok
Project to an original full tuple $i$ of length $j+j'+2$. Each
native coface projects to the original tuple with its indicated
entry deleted and restricts from that intersection to the full
intersection. The original differential is its finite signed sum.
Selecting a subtuple and restricting to a further original
intersection composes the same restriction maps. Consequently,
every projected coface of the product is the original exterior
product after the two corresponding composite subtuple selections.

Split the deletions into positions $0,\ldots,j$ and
$j+1,\ldots,j+j'+1$. For $k\leq j$, deleting $k$ before selecting
the prefix is the same as selecting the prefix of length $j+2$
and deleting its $k$th entry; the suffix now starts in position
$j+1$. These are exactly the source-left cofaces except their last
face. For $k=j+1+l$, the prefix stays unchanged and the suffix
has its $(l+1)$st entry deleted. These are exactly the source-right
cofaces except their first face. The latter signs agree because
$(-1)^{j+1+l}=(-1)^j(-1)^{l+1}$.

The omitted last-left and first-right faces have the identical
prefix ending in position $j$ and identical suffix beginning in
position $j+1$. Their entire original factor maps therefore coincide.
Their coefficients $(-1)^{j+1}$ and $(-1)^j$ sum to zero.
This proves the finite alternating-sum identity in the original
additive group of morphisms to this original tuple factor.
The actual tensor construction is additive in each original input
map, so its signed sums are precisely the tensor maps of the
original source differentials. Thus the projected identity is
exactly the displayed product rule. All original product projections
jointly detect morphism equality, proving the whole sheaf-map
identity.
\end{proof}
